% Section 7 — Theorem 1 as a Design-Space Map (~0.9 pages, weight 8%)
% Closing claim: every surveyed mitigation instantiates (i), (ii), or A4-only.
%
% REVIEW STATUS: Draft. Table rows to be cross-checked against §6 Obs 1 for
% consistency. Espresso row updated after architecture.md manual read (TF-3 PASS).
% T3rn added as new capability (ii) positive example.
% Appendix A boundary case needed for T3rn (sequential proof relay + rollback).
% (preconfer selection) and Task 2.1 (architecture analysis).

\section{Theorem 1 as a Design-Space Map}
\label{sec:designspace}

Theorem~\ref{thm:a2-necessary} partitions the design space of cross-rollup
sequencing mechanisms into three exhaustive classes: those providing capability
(i), those providing capability (ii), and those providing only A4 guarantees.
This section validates that every deployed or actively proposed mitigation
falls into exactly one of these classes, confirming the theorem is non-vacuous
and covers the actual design landscape.

Table~\ref{tab:designspace} maps 12 deployed and proposed mechanisms across
the full classification. Boundary cases — where the classification requires
careful argument — are analyzed in Appendix~\ref{app:boundary}; we summarize
the classification rationale here.

\begin{table*}[t]
\centering
\small
\caption{Design-space map: every surveyed cross-rollup mechanism instantiates
  Theorem~\ref{thm:a2-necessary} capability (i), (ii), (i)+(ii), or A4 only.
  ``Closes A2'' indicates whether the mechanism, if correctly deployed,
  provides Definition~\ref{def:a2} execution atomicity.
  $\dagger$ Boundary case; see Appendix~\ref{app:boundary}.}
\label{tab:designspace}
\renewcommand{\arraystretch}{1.15}
\begin{tabular}{@{}llllp{5.2cm}@{}}
\toprule
Mechanism & Class & Closes A2? & Evidence type & Classification rationale \\
\midrule
%--- Capability (i): Pre-execution evidence ---
ZK proof gating
  & (i) & Yes & Cryptographic
  & Commitment requires verifiable proof that both STFs succeed;
    $\phi(B)=\mathsf{true}$ is the ZK proof. \\
Eager simulation (sequencer executes)
  & (i) & Yes (local)
  & Operational
  & Sequencer runs $\delta_1, \delta_2$ locally before committing;
    not lazy by Definition~\ref{def:lazy}. \\
Espresso HotShot (confirmation layer)$^\dagger$
  & A0/A1 & No (lazy)
  & Documentation
  & \emph{``L2 validators receive blocks and execute the STFs for their
    rollups''} after HotShot confirmation~\cite{Espresso-Architecture}.
    Commits to ordering; no execution outcome in commitment. \\
Espresso + Hyperlane (Presto pattern)
  & (ii) & Yes (relay)
  & Contractual + code
  & \texttt{onlyMailbox} modifier on destination: destination executes
    only after Hyperlane relay delivers source event~\cite{Espresso-Presto}. \\
\midrule
%--- Capability (ii): Partial-effect prevention ---
HTLC (hash time-lock contract)
  & (ii) & Yes & Cryptographic
  & Hashlock ensures $\tau_2$ can only execute if $\tau_1$'s preimage
    is revealed; partial finality blocked at protocol level. \\
Across Protocol (UMA + escrow)$^\dagger$
  & (ii) & Yes (for fills)
  & Contractual
  & Filler escrow released only upon UMA oracle confirmation of both
    legs; revert path preserved until confirmation. \\
Astria application-level escrow
  & (ii) & Yes (app-level)
  & Contractual
  & Application holds funds in escrow over sequencer bundle; reverts
    all if either leg fails within the rollup timeout. \\
UniswapX cross-chain (signed orders)$^\dagger$
  & (ii)+A4 & Partial
  & Contractual + economic
  & Fill is protected by signed-order constraints (ii); residual
    timing gap covered by filler economics (A4). \\
t3rn (executor + proof relay)
  & (ii) & Yes (rollback)
  & Contractual + cryptographic
  & Executors submit proofs to settlement layer; protocol rolls back
    all legs on any failure. \emph{``No partial outcomes, no stranded
    assets, just a clean reset.''}~\cite{t3rn-docs} \\
\midrule
%--- Hybrid (i)+(ii) ---
Shutter Network (threshold encryption)$^\dagger$
  & (i)+(ii) & Yes (honest threshold)
  & Cryptographic
  & Encrypted intent prevents selective exercise of commitment (i);
    threshold decryption prevents partial reveal (ii). \\
\midrule
%--- A4 only: compensation, not prevention ---
Bond / slashing (generic)
  & A4 only & No & Economic
  & Bond compensates for A2 failure; does not provide $\phi(B)$
    (no (i)) and does not revert partial effects (no (ii)). \\
EigenLayer restaking (AVS preconf)$^\dagger$
  & A4 only & No & Economic
  & Restaking extends slashing to additional stake; same execution
    model as generic bond. See Appendix~\ref{app:boundary}. \\
MEV-Boost relay attestation
  & A4 only & No & Economic
  & Relay attests to block header inclusion, not execution outcome;
    no success predicate, no revert authority. \\
mev-commit / Bolt preconfirmation$^\dagger$
  & A4 only & No & Economic
  & Signed preconf commits to inclusion; bond slashed on
    violation. No STF execution in commitment path. \\
\bottomrule
\end{tabular}
\end{table*}

\paragraph{Completeness of the classification.}
Every mechanism in Table~\ref{tab:designspace} falls into one of the three
classes. The hybrid (i)+(ii) class (Shutter) is a special case where both
capabilities are simultaneously provided; Theorem~\ref{thm:a2-necessary}
requires \emph{at least one}, and the hybrid exceeds this requirement.

\paragraph{A4-only mechanisms cannot be upgraded to A2.}
A common design pattern is to combine multiple A4-only mechanisms (e.g.,
EigenLayer restaking + relay attestation + larger bond) in the expectation
that their combination approaches A2. This intuition is incorrect: the sum
of A4 guarantees is still A4. None of the A4-only mechanisms in
Table~\ref{tab:designspace} provides a success predicate or a revert
authority; their combination does not create either capability. The design
escape route is not to combine A4 mechanisms but to introduce capability (i)
or (ii).

\paragraph{Positive-finding candidates.}
Three mechanisms in Table~\ref{tab:designspace} provide A2 or conditional A2:
ZK proof gating, HTLC, and Across Protocol. These are the architectures that
correctly instantiate Theorem~\ref{thm:a2-necessary} and should serve as
templates for new designs seeking to close the A2 gap. The constraint for the
candidate analysis in \S\ref{sec:pefo} is precisely that the target application
does \emph{not} already instantiate these patterns (AX-4 filter in
\S\ref{sec:candidate-selection}).

\paragraph{Design implication.}
An application targeting A2 atomicity has two viable paths: embed a ZK proof
of joint execution outcome into the sequencer commitment (capability (i), high
engineering cost, strong guarantee), or implement application-level escrow
with revert authority (capability (ii), lower cost, requires rollup cooperation
for the revert path). Hybrid approaches (Shutter) provide stronger guarantees
at the cost of liveness dependency on the threshold committee.
